% !PRINTED 416
\weberpage{416}
%<<CONT
besteht, die für alle den Relationen (1) genügenden Werthe der $x$
befriedigt ist und worin $z$, immer mit Rücksicht auf die Relationen (1),
ungeändert bleibt, wenn die $x$ den Permutationen der Gruppe
$\mathfrak{A}$ unterworfen werden. Die Function $u$ soll dadurch in
$\nu$ verschiedene Functionen
\weq{3}{u, u_1, \ldots, u_{\nu-1}}
übergehen, so dass die Grössen (3) die Wurzeln der Gleichung (2) sind, und
wenn $\nu > 1$ ist, (2) eine Resolvente der Gleichung für die $x$ ist.

In dieser Allgemeinheit haben wir für die Lösung der Frage bis jetzt noch
keinen Angriffspunkt. Wir fügen daher eine weitere beschränkende
Voraussetzung hinzu, nämlich: \gesperrt{Es soll möglich sein, für die $x_h$
rationale Functionen einer Variablen $t$}
\weq{4}{x_h = \varphi_h(t)}
\gesperrt{zu setzen, so dass die Relationen (1) identisch befriedigt sind, und dass
zugleich $z$ nicht von $t$ unabhängig
wird}\footnote{Betrachtet man die $x$ und die $u$ als Coordinaten je eines
Punktes $X$ und $U$ in einem Raume von $n$ und $\nu$ Dimensionen, so
bestimmen die Relationen (1) für $x$ eine Mannigfaltigkeit von weniger
als $n$ Dimensionen. Jedem Punkte $X$ entspricht ein Punkt $U$, und der
Gesammtheit der $X$ nach (2) eine Mannigfaltigkeit der $U$ von einer
Dimension (einer Curve). Unsere Voraussetzung ist die, dass dies eine
sogenannte rationale oder unicursale Curve oder (in der Sprache der
Functionentheorie) eine Curve vom Geschlecht Null sei.}.

Um gleich hier ein Beispiel anzuführen, mit dem wir uns später noch
eingehend beschäftigen werden, erwähnen wir den Fall von fünf Variablen
$x_0, x_1, x_2, x_3, x_4$, die den Bedingungen
\weq{5}{\Sigma x = 0, \quad \Sigma x^2 = 0}
genügen, und die also die Wurzeln einer \gesperrt{Hauptgleichung}
5\textsuperscript{ten} Grades sind (Bd. I, \S.\,54). Dass die
Voraussetzung, die wir gemacht haben, in diesem Falle zutrifft, zeigt die
Darstellung in Bd. I, \S.\,74, (4):
\[
y = t_3 F_0 + t_2(\alpha_2 F_1 + \alpha_1 F_2 + \alpha_0 F_3),
\]
was, wenn die dort angeführten Bestimmungen getroffen sind, für die $x$
gesetzt, den Bedingungen (5) für alle Werthe des Verhältnisses
$t = t_2 : t_3$ genügt.
